%!TEX root = ../thesis.tex
%*******************************************************************************
%****************************** Third Chapter **********************************
%*******************************************************************************
\chapter[SenseMyStreet]{SenseMyStreet - Sensor Commissioning Toolkit for Communities}

% **************************** Define Graphics Path **************************
\ifpdf
    \graphicspath{{Chapter3/Figs/Raster/}{Chapter3/Figs/PDF/}{Chapter3/Figs/}}
\else
    \graphicspath{{Chapter3/Figs/Vector/}{Chapter3/Figs/}}
\fi

Describes the creation and implementation of the sensor commissioning toolkit and evaluation through use cases in the field.

\section{Introduction}

\section{Background/Related work}
Participation in smart city

Data-driven

Participatory sensing, VGI

Citizen science (Hakalay) from crowdsourcing to extreme citizen science (group projects)

Pittsburgh air pollution monitoring, aircast

Tools?? Maybe don’t focus on tools because they are just tools.

Community organising

Community networks - (John M Carroll and Mary Beth Rosson. 2003. A Trajectory for Community Networks. The Information Society 19: 381–393.)

\section{The process/Design of the sensor commissioning toolkit/Designing sensor commissioning toolkit}
Workshop in Gosforth (Beyond pins on the map)

Planning proposals and changes (Streets for People, Cycle ambition fund)

Activist groups informing (Groups of concerned citizens) 

Community support

Press interest

SPAN

Council/Government work

Available Resources

¥	Identifying the need for data (Discovery Phase). Through conducting interviews and observations of existing groups map out how they are using/not using data. 

¥	Making data and technology accessible (Design/Develop Phase). Undertaking a co-design process to make platforms and tools that respond to the needs of citizens. 
Starting with paper based methods and then moving to digital.

¥	Meaningful presentation of data (Deploy Phase). Ascertain what visualisations and formats of data are more accessible and useful for the sensemaking processes in different contexts. 

¥	Taking collective action by purposing data (Sensemaking Phase). Supporting the processes of sensemaking and assessing the use of data for decision-making and developing new solutions. 

\section{SenseMyStreet Toolkit}
Describe the process step by step